\begin{textsample}{NEG Dim 4 – 2021\_c6 – Score 1.00 – 2021\_c6/enade\_2021\_tec\_red\_comp\_4311331117.txt}  \label{ex:f4_neg_289}
CAPÍTULO 6 ANÁLISE TÉCNICA DA PROVA Este capítulo tem por objetivo apresentar o desempenho dos estudantes concluintes de Tecnologia em Redes de Computadores no Enade/2021.
Para isso, foram calculadas as estatísticas básicas da prova em seu todo ( seção 6.1.1 ), bem como as estatísticas relacionadas aos Componentes de Formação Geral ( seção 6.1.2 ) e de Conhecimento Específico da Área ( seção 6.1.3 ).
Dadas as suas características, foram analisadas, em separado, as questões objetivas ( seção 6.2 ) e as questões discursivas ( seção 6.3 ).
Tomando -se como base as duas Questões Discursivas do Componente de Formação Geral, nas seções 6.3.1.6 a 6.3.1.8, são apresentados os resultados e os comentários sobre a correção das respostas em relação à Língua Portuguesa.
Para as questões objetivas, foram disponibilizados os Índices de Facilidade e de Discriminação Ponto-Bisserial, também em separado, para os Componentes de Formação Geral ( seção 6.2.1 ) e de Conhecimento Específico ( seção 6.2.2 ).
De cada componente, uma das questões foi escolhida para exemplificar a análise gráfica, relacionando -se as alternativas escolhidas pelos estudantes ( inclusive o gabarito ) com o número de acertos no componente.
No Anexo I, apresenta -se a íntegra da análise gráfica para todas as questões objetivas.
Para cada uma das questões discursivas, os conteúdos dos tipos mais comuns de respostas dos estudantes são apresentados e comparados com o padrão de resposta esperado ( ver Anexo IX com o padrão de respostas ).
Nas tabelas, constam as seguintes estatísticas das notas19 : média do desempenho na prova, erro padrão da média, desvio padrão, nota mínima, mediana e nota máxima para cada um de seus componentes.
Tais estatísticas contemplam o total de estudantes concluintes da Área de Tecnologia em Redes de Computadores inscritos e presentes à prova do Enade/2021, tendo em vista agregações, ou por Grandes Regiões e o país como um todo, ou por Categoria Administrativa, Organização Acadêmica da IES e Modalidade de Ensino.
Em relação aos gráficos de distribuição de notas, o intervalo considerado foi de 10 unidades, aberto à esquerda e fechado à direita, com exceção do primeiro intervalo, [ 0 ; 10 ], fechado em ambos os extremos.
Para os gráficos de distribuição das notas das questões discursivas, foram consideradas mais duas categorias : questão em branco20 e nota zero. 19 Uma definição dessas estatísticas pode ser encontrada no Glossário. 20 Nesse grupo estão incluídas também as respostas classificadas como nulas ou desconsideradas. 6.1 ESTATÍSTICAS BÁSICAS DA PROVA Esta seção apresenta estatísticas selecionadas e histogramas da nota geral ( 6.1.1 ) e de cada componente : Formação Geral ( 6.1.2 ) e Conhecimento Específico ( 6.1.3 ).
São também apresentadas estatísticas selecionadas de subpopulações, caracterizadas por Grande Região, Categoria Administrativa, Organização Acadêmica e Modalidade de Ensino. 6.1.1 ESTATÍSTICAS BÁSICAS GERAIS Na Tabela 6.1, são apresentadas as Estatísticas Básicas da prova, por Grande Região, dos estudantes concluintes de Tecnologia em Redes de Computadores.
A Média das notas da prova, como um todo ( nas seções seguintes serão analisados os Componentes de Formação Geral e de Conhecimento Específico ), foi 28,5, sendo que os estudantes da região Centro-Oeste obtiveram a Média mais baixa ( 25,3 ), e os da região Sul obtiveram a Média mais alta ( 33,7 ).
As demais Médias foram : 28,2, na região Norte, 31,8, na região Nordeste e 26,7, na região Sudeste.
O Desvio padrão para o Brasil, como um todo, foi 11,8, sendo o maior Desvio padrão encontrado na região Nordeste ( 12,9 ), e o menor, na região Centro-Oeste ( 9,4 ), indicando uma dispersão um pouco menor das notas dessa última região.
A região que obteve a maior nota Máxima foi a Nordeste ( 67,5 ), ao passo que a região que atingiu a menor nota Máxima foi a Centro-Oeste ( 51,9 ).
A Mediana do Brasil, como um todo, foi 26,6, sendo a maior Mediana obtida na região Sul ( 32,2 ), e a menor obtida na região Centro-Oeste ( 24,3 ).
A nota Mínima do Brasil foi obtida na região Sudeste, onde foi zero.
A maior nota Mínima, entre as cinco Grandes Regiões, foi 7,2, na região Sul, as demais regiões obtiveram as seguintes notas Mínima : 5,3 ( Norte ), 3,4 ( Nordeste ) e 7,1 ( Centro-Oeste ).
Considerando -se as notas segundo Grande Região, observa -se que existe diferença estatisticamente significativa ao nível de 95 \% entre as Médias obtidas na região Sul e nas regiões Norte, Sudeste e Centro-Oeste, e entre a região Nordeste e as regiões Sudeste e Centro-Oeste.
O comportamento das notas dos estudantes de todo o Brasil pode ser observado no Gráfico 6.1, que apresenta um histograma com a distribuição das mesmas : uma distribuição unimodal com a moda no intervalo ( 20 ; 30 ].
Na Tabela 6.2, são apresentadas informações referentes à Média da nota final, desagregadas por Categoria Administrativa, por Organização Acadêmica e por Modalidade de Ensino.
A Média das notas da prova, como um todo, foi 28,5.
Em relação a Categoria Administrativa, os estudantes das IES Públicas obtiveram Média mais alta ( 33,2 ) que a Média nacional, e os das IES Privadas obtiveram Média mais baixa ( 27,2 ).
Observa -se que existe diferença estatisticamente significativa entre as Médias das notas das IES Públicas e as das Privadas.
A diferença entre as Médias das regiões Sul e Centro-Oeste ( 8,4 ), a maior e a menor Média, é superior à diferença entre a das IES Públicas e Privadas ( 6,0 ), não chegando a caracterizar uma maior diversidade administrativa do que regional.
O Desvio padrão para as IES Públicas ( 12,7 ) foi superior ao do Brasil, como um todo ( 11,8 ), indicando uma dispersão um pouco maior das notas nesta Categoria Administrativa.
No tocante à Organização Acadêmica, as Faculdades ( 29,2 ) e os CEFET/IF ( 35,5 ) obtiveram Média mais alta que a nacional ( 28,5 ).
A Média das Universidades ( 27,0 ) e dos Centros Universitários ( 28,0 ) foi menor do que a nacional.
Constata -se que existe diferença estatisticamente significativa, ao nível de 95 \%, entre as Médias dos CEFET/IF em relação às demais Organizações Acadêmicas, assim como entre as Universidades e Faculdades.
A Média da Modalidade Educação a Distância ( 27,1 ) foi inferior à Média da Educação Presencial ( 29,6 ), e a diferença entre elas é estatisticamente significativa, ao nível de 95 \%. 6.1.2 ESTATÍSTICAS BÁSICAS NO COMPONENTE DE FORMAÇÃO GERAL Na Tabela 6.3, são apresentadas as Estatísticas Básicas em relação ao componente da prova que avalia a Formação Geral dos estudantes concluintes.
A nota de Formação Geral é obtida levando -se em conta os acertos nas oito questões objetivas e a nota média obtida nas Questões Discursivas 1 e 2, sendo a nota de cada questão discursiva obtida levando -se em conta 80 \% da nota relativa ao conteúdo e 20 \% da nota relativa ao desempenho linguístico.
Os estudantes de todo o Brasil obtiveram desempenho médio de 28,7.
Quanto à variabilidade, o Desvio padrão das notas dos estudantes do Brasil foi 14,0.
A menor Média foi obtida na região Centro-Oeste ( 27,0 ) e a maior, na região Sul ( 31,0 ).
As demais Médias foram : 28,5, na região Norte, 30,8, na região Nordeste e 27,7, na região Sudeste.
Já o maior Desvio padrão foi obtido na região Norte ( 14,7 ), e o menor, na região Sul ( 13,4 ).
Os demais desvios foram : 13,9, na região Nordeste, 14,1, na região Sudeste e 13,5, na região Centro-Oeste.
A maior nota no Componente de Formação Geral da prova do Enade/2021 foi 76,6, obtida por, pelo menos, um estudante na região Sudeste.
A menor nota Máxima foi obtida na região Norte ( 57,8 ).
A Mediana do Brasil, como um todo, foi 28,1, sendo a menor Mediana encontrada na região Norte ( 26,1 ), e a maior, nas regiões Nordeste e Sudeste ( 30,0 ).
A nota Mínima nessa parte foi zero em quase todas as regiões, exceto na região Sul ( 2,0 ).
Considerando -se as notas, segundo Grande Região, observa -se que existe diferença estatisticamente significativa entre as Médias das notas do Componente de Formação Geral obtidas na região Nordeste e nas regiões Sudeste e Centro-Oeste, assim como entre as regiões Sudeste e Sul.
No Gráfico 6.2, é possível verificar a avaliação do desempenho dos estudantes do Componente de Formação Geral, a partir do histograma da distribuição das notas correspondentes.
A distribuição é unimodal, com moda em ( 20 ; 30 ], o mesmo intervalo modal da distribuição de notas da prova, como um todo ( Gráfico 6.1 ).
Nota -se, ainda, que, no Gráfico 6.2, as notas apresentam maior dispersão do que as no Gráfico 6.1 ( distribuição das notas da prova ), confirmada pela comparação dos desvios padrões : 11,8 para a nota da prova, como um todo, e 14,0 para o Componente de Formação Geral.
Na Tabela 6.4, são apresentadas as informações referentes ao desempenho dos concluintes do Componente de Formação Geral, em diferentes agregações : Categoria Administrativa, Organização Acadêmica e Modalidade de Ensino.
Observa -se que existe diferença estatisticamente significativa entre as Médias dos tipos de Categoria Administrativa : a maior Média obtida por estudantes de IES Públicas ( 32,4 ) e a menor, pelos de IES Privadas ( 27,7 ).
Nota -se que, também, há diferença estatisticamente significativa entre as Médias dos CEFET/IF ( 33,7 ) em relação aos demais tipos de Organização Acadêmica : 27,9 para as Universidades, 27,7 para os Centros Universitários e 29,0 para as Faculdades.
No que se refere a Modalidade de Ensino, a Educação Presencial apresentou maior Média do que a Educação a Distância ( 29,2 e 28,0, respectivamente ).
Tais Médias não apresentam diferença estatisticamente significativa. 6.1.3 ESTATÍSTICAS BÁSICAS DO COMPONENTE DE CONHECIMENTO ESPECÍFICO Na Tabela 6.5, são apresentadas as Estatísticas Básicas referentes ao Componente de Conhecimento Específico da Área de Tecnologia em Redes de Computadores.
A nota do componente de Conhecimentos Específicos leva em conta as notas da parte objetiva e a média das notas obtidas nas três questões discursivas parte discursiva da Área.
A Média do desempenho dos estudantes do Brasil, como um todo, foi 28,4.
A maior Média foi obtida na região Sul ( 34,5 ), e a menor, na região Centro-Oeste ( 24,7 ).
As demais Médias foram : 28,0 na região Norte, 32,1 na região Nordeste e 26,3 na região Sudeste.
Quanto à variabilidade das notas, o Desvio padrão do Brasil, como um todo, foi 13,6, sendo o maior Desvio padrão observado na região Nordeste ( 15,0 ) e o menor, na região Centro-Oeste ( 10,9 ).
Os demais desvios foram : 12,2 na região Norte, 12,6 na região Sudeste e 14,9 na região Sul.
A Mediana das notas dos estudantes de todo o Brasil foi 26,9.
A maior Mediana ocorreu na região Sul ( 32,1 ), e a menor, na região Centro-Oeste ( 22,4 ).
As demais Medianas foram : 26,9 na região Norte, 30,3 na região Nordeste e 23,9 na região Sudeste.
A nota Máxima do Brasil, como um todo, foi 82,1, sendo obtida por pelo menos um estudante na região Nordeste.
As demais notas Máximas foram : 62,4 na região Norte, 74,8 na região Sudeste, 78,8 na região Sul e 56,7 na região Centro-Oeste.
A nota Mínima foi zero na região Sudeste.
As notas Mínimas das demais regiões foram 4,5.
Observa -se que existe diferença estatisticamente significativa entre a região Sul e as regiões Norte, Sudeste e Centro-Oeste na análise das Médias das notas do Componente de Conhecimento Específico, assim como entre a região Nordeste e as regiões Sudeste e Centro-Oeste.
Assim como nos Gráficos 6.1 e 6.2, no Gráfico 6.3, apresentado a seguir, é possível ser feita uma avaliação do desempenho de concluintes de Tecnologia em Redes de Computadores, em relação ao Componente de Conhecimento Específico, com um histograma da distribuição das notas correspondentes.
Esta também é uma distribuição unimodal, e o grupo modal é o ( 20 ; 30 ], o mesmo do grupo modal da prova, como um todo, e para a Formação Geral.
Na Tabela 6.6, é apresentada uma comparação dos resultados em relação a Categoria Administrativa, a Organização Acadêmica e a Modalidade de Ensino, agora levando -se em conta o desempenho de estudantes do Componente de Conhecimento Específico da prova.
No que se refere a Organização Acadêmica, a maior Média foi a dos CEFET/IF ( 36,2 ), vindo a seguir a das Faculdades ( 29,2 ), a dos Centros Universitários ( 28,0 ) e, por fim, a das Universidades ( 26,7 ).
O maior Desvio padrão, e igual ao valor para o Brasil, como um todo, foi o dos CEFET/IF ( 15,4 ).
Os CEFET/IF obtiveram a maior nota Máxima ( 82,1 ).
As Universidades obtiveram nota Máxima 74,8, os Centros Universitários, 67,1, e as Faculdades, 78,8.
As Medianas foram : 25,4 nas Universidades e nos Centros Universitários, 26,9 nas Faculdades e 31,3 nos CEFET/IF.
A nota Mínima foi zero para as Universidade, 4,5 para os Centros Universitários e as Faculdades, e 9,9 para os CEFET/IF.
Observa -se que não existe diferença estatisticamente significativa, ao nível de 95 \%, no Componente de Conhecimento Específico entre as notas dos Centros Universitários em relação às Universidade e Faculdades.
Quanto a Categoria Administrativa, observa -se um comportamento semelhante àquele da parte de Formação Geral e à prova, como um todo, ou seja, existe diferença estatisticamente significativa entre as Médias das IES Públicas ( 33,5 ) e as das IES Privadas ( 27,0 ).
Neste caso, também, a maior Média foi obtida por estudantes de IES Públicas de ensino.
Quanto ao comportamento das Médias, no que toca a Modalidade de Ensino, há diferença estatisticamente significativa entre a Média da Educação Presencial ( 29,7 ) e a da Educação a Distância ( 26,7 ). 6.2 ANÁLISE DAS QUESTÕES OBJETIVAS Esta seção apresenta estatísticas selecionadas e histogramas das Questões Objetivas de Formação Geral ( 6.2.1 ) e de Conhecimento Específico ( 6.2.2 ).
São, também, apresentadas e comparadas as médias das subpopulações caracterizadas por Grande Região. 21 6.2.1 COMPONENTE DE FORMAÇÃO GERAL Na Tabela 6.7, são apresentadas as Estatísticas Básicas relativas às oito Questões Objetivas do Componente da prova que abrange a Formação Geral dos estudantes.
A Média do Brasil foi 36,2.
A menor Média foi encontrada na região Centro-Oeste ( 34,7 ), e a maior, na região Sul ( 38,4 ).
As demais Médias foram : 35,0 na região Norte, 37,5 na região Nordeste e 35,6 na região Sudeste.
O Desvio padrão do Brasil foi 19,1, sendo o maior Desvio padrão encontrado na região Sudeste ( 19,5 ), e o menor Desvio padrão encontrado na região Norte ( 15,4 ).
As Medianas ( 37,5 ) e as notas Mínimas ( 0,0 ) foram iguais para todas as regiões.
A nota Máxima 100,0 foi alcançada na região Sudeste. 21 As fórmulas utilizadas para o coeficiente de correlação ponto-bisserial e para o índice de facilidade estão apresentadas no capítulo 1.
Na Tabela 6.8, são apresentados o Índice de Facilidade e o Índice de Discriminação ( Ponto-Bisserial ) para cada uma das Questões Objetivas do Componente de Formação Geral.
Quanto ao Índice de Facilidade, foram usadas as seguintes cores para diferenciar o nível de dificuldade da questão :  Azul para as questões classificadas com índice Muito fácil ( \&gt ; =0,86 ), verde para as questões classificadas com índice Fácil ( 0,61 a 0,85 ), amarelo para as questões classificadas com Médio ( 0,41 a 0,60 ), vermelho para as questões classificadas com Difícil ( 0,16 a 0,40 ) e roxo para as questões classificadas com Muito difícil ( \&lt ; =0,15 ).
Já quanto ao Índice de Discriminação, foram usadas as seguintes cores para qualificar a questão :  As questões classificadas com índice Fraco receberam a cor vermelha ( \&lt ; =0,19 ), as classificadas com Médio receberam a cor amarela ( 0,20 a 0,29 ), as classificadas com Bom receberam a cor verde ( 0,30 a 0,39 ) e as classificadas com Muito bom ( \&gt ; =0,40 ) receberam a cor azul.
As questões objetivas do Componente de Formação Geral, segundo o Índice de Facilidade, foram assim avaliadas : nenhuma teve o Índice de Facilidade classificado como Muito fácil ou como Fácil.
Quatro questões foram consideradas com índice de dificuldade Médio, situando -se no intervalo entre 0,41 e 0,60 do Índice de Facilidade, ou seja, houve entre 41,0 \% e 60,0 \% de acertos, enquanto outras quatro questões foram classificadas na categoria Difícil, apresentando acertos na faixa entre 16 \% e 40 \%.
Nenhuma questão apresentou menos de 16 \% de acertos, razão pela qual não houve questão classificada como Muito difícil.
O Índice de Facilidade variou de 0,18 a 0,54, e o de Discriminação, de 0,23 a 0,47.
Seis questões, quanto ao Índice de Discriminação, foram tidas como Muito bom, uma questão foi tida como Bom e uma questão foi tida como Médio.
Na Tabela 6.9, é apresentada a distribuição das questões, levando -se em conta, simultaneamente, a classificação dos dois índices.
Como já comentado, nenhuma questão foi anulada.
As seis questões que tiveram Índice de Discriminação Muito bom figuraram em dois níveis de dificuldade Difícil e Médio : duas classificadas na categoria Difícil ( Questões 1 e 4 ) do Índice de Facilidade, e quatro, na categoria Médio ( Questões 3, 5, 6 e 7 ).
Em particular, a Questão 7 foi a que apresentou o maior poder discriminatório, com índice 0,47, e foi considerada Médio em termos de Facilidade, com uma proporção de 0,54 acertos, o máximo de acertos alcançado.
A Questão 2 foi considerada Difícil, com Índice de Facilidade 0,18.
O Índice de Discriminação foi Bom para a Questão 2 ( 0,37 ).
Na Tabela 6.9, é apresentada a distribuição das questões, levando -se em conta, simultaneamente, a classificação dos dois índices.
Como já comentado, nenhuma questão foi anulada.
As seis questões que tiveram Índice de Discriminação Muito bom figuraram em dois níveis de dificuldade Difícil e Médio : duas classificadas na categoria Difícil ( Questões 1 e 4 ) do Índice de Facilidade, e quatro, na categoria Médio ( Questões 3, 5, 6 e 7 ).
Em particular, a Questão 7 foi a que apresentou o maior poder discriminatório, com índice 0,47, e foi considerada Médio em termos de Facilidade, com uma proporção de 0,54 acertos, o máximo de acertos alcançado.
A Questão 2 foi considerada Difícil, com Índice de Facilidade 0,18.
O Índice de Discriminação foi Bom para a Questão 2 ( 0,37 ).
No Gráfico 6.4, para exemplificar, analisa -se o comportamento da Questão de número 3 de Formação Geral.
Trata -se de uma questão considerada com nível Médio de dificuldade, em relação à Facilidade, e a que obteve o segundo maior Índice de Discriminação dessa parte da prova.
Neste gráfico, cada uma das cinco curvas representa o percentual de respostas em determinada alternativa da questão, em função do número de acertos dos estudantes nessa parte da prova ( Formação Geral/Múltipla Escolha ), antes de possíveis eliminações pelo Como nenhuma Questão Objetiva do Componente de Formação Geral foi anulada pela Comissão Assessora de Área o número de acertos alcança o máximo de oito questões para esse componente.
A curva em verde corresponde à alternativa B, a alternativa correta para essa questão.
Observa -se que entre os estudantes com menor número de acertos, nessa parte do exame, a situação mais frequente foi a escolha da alternativa E ( em vermelho ), incorreta.
Por exemplo, entre os estudantes que acertaram duas respostas, 51,8 \% escolheram a alternativa E ( em vermelho ), 28,4 \% escolheram a alternativa B ( em verde, gabarito ), 10,2 \% escolheram a alternativa D ( em laranja ), 5,4 \% escolheram a alternativa A ( em azul ) e 4,0 \% escolheram a alternativa C ( em preto ).
Dentre os que acertaram duas respostas entre as questões de múltipla escolha de Formação Geral, 0,2 \% deixou a resposta questão em branco e 0,0 \% marcou mais de uma alternativa, invalidando a questão.
À medida que o número de acertos aumenta, indicando desempenho melhor nessa parte da prova, aumenta, concomitantemente, a proporção de estudantes que selecionaram a alternativa correta B, atingindo 100 \% para os estudantes com oito acertos.
Essa análise permite verificar como a questão discriminou os grupos de desempenho, justificando -se o alto índice ( 0,46 ) obtido na questão.
Cumpre notar que não é possível inferir deste gráfico nem o Índice de Facilidade, que seria uma média da proporção ponderada pela quantidade de estudantes com cada uma das notas, nem o Índice de Discriminação Ponto-Bisserial, por razão equivalente.
No caso extremo, no qual a grande concentração dos acertos dos estudantes fosse abaixo de quatro, o Índice de Facilidade seria obrigatoriamente abaixo de 60,0 \% ( neste exemplo ).
Caso a concentração fosse em seis acertos ou mais, o índice seria obrigatoriamente acima de 80,0 \%.
Os gráficos relativos às demais questões de Formação Geral constam do Anexo I. 6.2.2 COMPONENTE DE CONHECIMENTO ESPECÍFICO Na Tabela 6.10, são apresentadas as Estatísticas Básicas em relação às Questões Objetivas do Componente de Conhecimento Específico da prova de Tecnologia em Redes de Computadores por Grande Região.
Nessa parte da prova, como será discutido ainda nesta seção, nenhuma Questão Objetiva do Componente de Conhecimentos Específicos foi anulada pela CAA e oito questões não foram usadas no cômputo das notas por terem sido descartadas pelo critério do Ponto-Bisserial.
Assim, as notas foram calculadas com base em 19 das 27 Questões Objetivas válidas de Conhecimentos Específicos.
A Média do Brasil deste componente foi 31,4.
A menor Média foi observada na região Centro-Oeste ( 27,6 ), e a maior, na região Sul ( 38,0 ).
O Desvio padrão de todo o Brasil foi 14,9, sendo o menor Desvio padrão encontrado na região Centro-Oeste ( 12,4 ), e o maior, na região Nordeste ( 16,4 ).
A Mediana de todo o Brasil foi 26,3, o mesmo valor da Mediana encontrada nas regiões Sudeste e Centro-Oeste.
Nas demais regiões o valor foi maior : Norte ( 29,0 ), Nordeste ( 31,6 ) e Sul ( 36,8 ).
A nota Máxima da prova foi 89,5, obtida nas Questões Objetivas do Componente de Conhecimento Específico, por, pelo menos, um estudante da região Nordeste.
Nas demais regiões, o valor foi : 68,4 na região Norte, 78,9 na região Sudeste, 84,2 na região Sul e 63,2 na região Centro-Oeste.
A nota Mínima foi zero, obtida por pelo menos um estudante na região Sudeste nas Questões Objetivas do Componente de Conhecimento Específico.
A nota Mínima nas demais regiões foi 5,3.
Na Tabela 6.11, são apresentados os Índices de Facilidade e de Discriminação ( Ponto-Bisserial ) das Questões Objetivas do Componente de Conhecimento Específico, para os estudantes de Tecnologia em Redes de Computadores.
Para facilitar a diferenciação das questões, serão usadas as mesmas cores da Tabela 6.8 para as diferentes classificações dos Índices de Facilidade e de Discriminação.
Observa -se que nenhuma questão foi anulada pela CAA, como já mencionado.
Sendo assim, as análises a seguir dizem respeito às 27 questões válidas.
A partir do Índice de Facilidade obtido, pode -se concluir que mais de quatro quintos das questões objetivas da prova foram consideradas Difícil ou Muito difícil : das 27 questões válidas, 21 foram classificadas como Difícil e duas como Muito difícil.
Uma questão foi classificada como Fácil, e outras três consideradas como Médio.
Nenhuma questão foi classificada como Muito fácil.
Já quanto aos Índices de Discriminação das Questões Objetivas do Componente de Conhecimento Específico da prova, obtém -se como resultado a seguinte classificação : duas das 27 questões válidas foram consideradas boas, enquanto oito delas tiveram Índice de Discriminação Muito bom.
Assim, para dez em 27 questões, os Índices de Discriminação foram Bom ou Muito bom.
Dentre as demais, nove delas foram classificadas como Médio, e outras oito, como Fraco, sendo 17, por conseguinte, a quantidade de questões nos dois patamares mais baixos de Discriminação.
Constata -se, assim, que a prova – no que se refere ao Componente de Conhecimento Específico – possuía média capacidade de discriminar entre aqueles que dominam ou não o conteúdo.
O Índice de Facilidade variou de 0,09 a 0,62 e o de Discriminação, de -0,10 a 0,47.
Na Tabela 6.12, é apresentada a distribuição das questões levando -se em conta, simultaneamente, a classificação dos dois índices.
Dentre as questões que alcançaram os maiores Índices de Discriminação, uma delas foi classificada com Índice de Discriminação Muito bom, a de número 24.
Seu índice foi 0,47, e, quanto ao Índice de Facilidade, foi classificada como Difícil, com 34 \% dos estudantes marcando a opção correta.
A Questão de número 28 foi a mais difícil entre as 27 questões específicas válidas, com baixo Índice de Facilidade, apenas 9 \% de acertos.
Essa questão apresentou poder discriminatório muito baixo, 0,02, o que comprova ter sido a mais difícil para os estudantes.
Destaca -se, também, a Questão 29, com Índice de Facilidade 0,11, o que, em termos percentuais, corresponde a 11 \% de estudantes que responderam acertadamente, sendo, também, 0,07 seu Índice de Discriminação.
Tais questões foram, portanto, pelo critério Ponto-Bisserial, consideradas inadequadas.
Por isso, as Questões 28 e 29, além das Questões 13, 19, 31, 32, 34 e 35, foram eliminadas do cômputo da nota final.
A título de exemplo das análises do comportamento das questões objetivas, no Gráfico 6.5, analisa -se a Questão 24 do Componente de Conhecimento Específico.
Essa questão foi considerada, pelas respostas dos estudantes avaliados na prova, como uma questão com nível de Facilidade Difícil, 0,34, ou seja, 34 \% dos estudantes assinalaram, acertadamente, a opção C correspondente ao gabarito.
Como já comentado, seu Índice de Discriminação foi igual a 0,47, classificado como Muito bom, o maior valor de Discriminação.
Neste gráfico, cada uma das cinco curvas representa o percentual de respostas em determinada alternativa da Questão 24, em função do número de acertos dos estudantes nessa parte da prova, antes de possíveis eliminações de questões pelo critério Ponto-Bisserial.
A alternativa correta C, representada no gráfico pela curva em preto, foi escolhida em maiores proporções pelos estudantes com desempenho melhor nessa parte da prova.
Já as alternativas incorretas, também denominadas distratores, foram selecionadas, principalmente, por aqueles com notas mais baixas.
No caso de Tecnologia em Redes de Computadores, como se observa no eixo horizontal do Gráfico 6.5, nenhum estudante acertou mais do que 19 questões dentre as 27 questões objetivas válidas, considerando -se o gabarito original.
A eliminação das oito questões com Índice de Discriminação Fraco do cômputo da nota não possibilitou a existência de nota Máxima igual a 100,0 por pelo menos um estudante ( visto na Tabela 6.10 ).
Observa -se que a soma não é 100 \%, por causa das questões não respondidas ou com mais de uma opção marcada.
Aqueles com um acerto, 50 \% marcaram a opção E ( linha vermelha ), 32,5 \% a opção D ( linha laranja ) e 17,5 \% a opção B ( linha verde ).
A proporção de estudantes que selecionaram a resposta correta C aumenta gradativamente, chegando a atingir 100 \% para 17 acertos ou mais, enquanto a proporção dos que escolheram alternativas

% matched lemmas: 
\end{textsample}
